% Replace subsection A with the following.
\subsection{Benchmark Construction}
We describe the Tanner graphs used in
Tables~\ref{tab:exact-exp} and~\ref{tab:hybrid-exp}.
Small degrees reduce the alternating branching factors in $\omega$,
motivating small-degree benchmarks. These are synthetic sparse
parity-check matrices, not standardized codes or decoding-performance tests.

Bases A--C use point--line incidence $y=sx+b$ over $\mathbb F_Q$,
the finite field with $Q$ elements, with
$x\in\{0,\ldots,d_c-1\}$ and $s\in\{0,\ldots,d_v-1\}$.
Their parameters $(Q,d_v,d_c)$ are $(7,3,6)$, $(5,3,4)$, and
$(11,4,8)$, respectively. Base D splits each degree-six point of the
symplectic generalized quadrangle over $\mathbb F_5$ into two
 degree-three variables, giving 312 variables, 156 checks, and girth eight.
A--D labels specify the base, lift size, and cyclic ($c$) or random ($r$) lift.
R1200 uses random socket pairing with degree-preserving switches removing
$4$-cycles. P1200 and P1600 use our degree-constrained PEG
variant~\cite{hu2005}; R/P labels give the number of variables.

We also include column-weight-two cycle-code benchmarks~\cite{malema2007}.
Pet64 and HS128 are cyclic lifts of the incidence Tanner graphs of the
Petersen and Hoffman--Singleton graphs. F3c32 and F5c64 similarly use
projective-plane incidence graphs over $\mathbb F_3$ and $\mathbb F_5$.
Ordinary edges become variables and vertices become checks, giving
$d_v=2$ and doubling girth to 10 or 12; label suffixes give lift sizes.
Generators, lift definitions, fixed seeds, and edge lists are in the
repository linked above. Every graph was checked for simplicity,
connectivity, degrees, and girth. Exact girth-cycle counts agree with
independent exhaustive NB enumeration and our KB implementation
on the same graphs.

% Replace Table I with the following.
\begin{table*}[t]
\centering
\caption{Same-target exact comparison; milliseconds, medians of seven runs.
The count tuple lists every even length from $g$ through $T$.
Speedup is KB time divided by our exact time; values below one favor KB.}
\label{tab:exact-exp}
\footnotesize
\setlength{\tabcolsep}{5pt}
\begin{tabular}{lccrrlrrr}
\hline
Graph & $(d_v,d_c)$ & $g$ & $n$ & $T$ & $(N_g,\ldots,N_T)$ & Ours & KB & Speedup\\
\hline
R1200 & $(4,8)$ & 6 & $1,200$ & 6 & $(1,564)$ & $4.080$ & $47.842$ & $11.73\times$\\
P1600 & $(3,4)$ & 10 & $1,600$ & 10 & $(11)$ & $5.318$ & $76.137$ & $14.32\times$\\
Pet64 & $(2,3)$ & 10 & $960$ & 10 & $(512)$ & $0.263$ & $0.610$ & $2.32\times$\\
HS128 & $(2,7)$ & 10 & $22,400$ & 10 & $(156,672)$ & $112.507$ & $263.448$ & $2.34\times$\\
F3c32 & $(2,4)$ & 12 & $1,664$ & 12 & $(6,624)$ & $1.825$ & $4.919$ & $2.70\times$\\
F5c64 & $(2,6)$ & 12 & $11,904$ & 12 & $(240,000)$ & $74.173$ & $167.401$ & $2.26\times$\\
A128r & $(3,6)$ & 6 & $5,376$ & 10 & $(176, 1,709, 8,987)$ & $561.784$ & $1466.804$ & $2.61\times$\\
A128c & $(3,6)$ & 6 & $5,376$ & 10 & $(20,480, 203,904, 1,069,056)$ & $466.414$ & $38.622$ & $0.08\times$\\
D1c & $(3,6)$ & 8 & $312$ & 14 & $(3,638, 15,670, 142,582, 1,111,500)$ & $983.688$ & $13.715$ & $0.01\times$\\
\hline
\end{tabular}
\end{table*}

% Replace Table II with the following.
\begin{table}[!ht]
\centering
\caption{Hybrid estimation over 100 runs.
Here $n$ is the number of variable nodes, $N_g$ is the exact girth-cycle
count, and $\lambda=gN_g/B$ is the expected number of closures in the
$n$-trial pilot.
Speedup is the median exact time over seven runs divided by the mean
hybrid time over 100 runs; values above one favor the hybrid method.
Hybrid time includes screening and exact fallback.
MARE is the mean of $100|\widetilde N_g-N_g|/N_g$ over all 100 outputs,
including exact returns.}
\label{tab:hybrid-exp}
\footnotesize
\setlength{\tabcolsep}{3.6pt}
\begin{tabular}{lccrrrrr}
\hline
Graph & $(d_v,d_c)$ & $g$ & $n$ & $N_g$ & $\lambda$
& Speedup & mare (\%)\\
\hline
A16c & $(3,6)$ & 6 & $672$ & $2,560$ & $10.24$ & $0.56\times$ & $2.9$\\
A128r & $(3,6)$ & 6 & $5,376$ & $176$ & $0.70$ & $0.89\times$ & $0.0$\\
A128c & $(3,6)$ & 6 & $5,376$ & $20,480$ & $81.92$ & $3.62\times$ & $7.3$\\
B128c & $(3,4)$ & 6 & $2,560$ & $4,608$ & $85.33$ & $1.61\times$ & $7.3$\\
C128c & $(4,8)$ & 6 & $11,264$ & $204,928$ & $99.58$ & $10.84\times$ & $8.2$\\
D16c & $(3,6)$ & 8 & $4,992$ & $57,680$ & $30.76$ & $4.34\times$ & $7.7$\\
R1200 & $(4,8)$ & 6 & $1,200$ & $1,564$ & $0.76$ & $0.96\times$ & $0.0$\\
P1200 & $(3,6)$ & 8 & $1,200$ & $65$ & $0.035$ & $0.90\times$ & $0.0$\\
P1600 & $(3,4)$ & 10 & $1,600$ & $11$ & $0.009$ & $0.92\times$ & $0.0$\\
Pet64 & $(2,3)$ & 10 & $960$ & $512$ & $80.00$ & $0.82\times$ & $7.9$\\
HS128 & $(2,7)$ & 10 & $22,400$ & $156,672$ & $100.74$ & $7.07\times$ & $7.6$\\
F3c32 & $(2,4)$ & 12 & $1,664$ & $6,624$ & $54.52$ & $2.04\times$ & $6.9$\\
F5c64 & $(2,6)$ & 12 & $11,904$ & $240,000$ & $92.16$ & $8.47\times$ & $7.6$\\
\hline
\end{tabular}
\end{table}

% Replace the paragraph after Table II with the following.
Table~\ref{tab:hybrid-exp} uses $r=100$ and $K=10n$.
Zero MARE results from exact fallback in every run; A16c has mixed routing.
All four cycle-code benchmarks use inverse sampling in every run,
including the girth-10 and girth-12 instances.
HS128, F3c32, and F5c64 favor hybrid estimation, whereas Pet64 shows
that high closure abundance alone need not yield a speedup when exact
counting is already inexpensive. The results cover sparse, mixed,
and cycle-rich regimes without implying superiority across LDPC families.

% Add these bibliography entries if not already present.
\bibitem{hu2005}
X.-Y. Hu, E. Eleftheriou, and D. M. Arnold,
``Regular and irregular progressive edge-growth Tanner graphs,''
\emph{IEEE Trans. Inf. Theory}, vol. 51, no. 1, pp. 386--398, 2005.

\bibitem{malema2007}
G. Malema and M. Liebelt,
``High girth column-weight-two LDPC codes based on distance graphs,''
\emph{EURASIP J. Wireless Commun. Netw.}, vol. 2007,
Art. no. 048158, 2007, doi: 10.1155/2007/48158.
